% ============================================================================
\section{Introduction}
\label{sec:intro}
% ============================================================================

Multi-teacher on-policy distillation (MOPD) merges several specialist teachers
into one student: the student generates a response, and the specialist for its
task scores that response token by token
\citep{ma2026mopdmultiteacheronpolicydistillation}. Existing recipes must also
choose a task mixture. Token-share balancing changes that mixture to compensate
for response-length imbalance \citep{gao2026openmopddiagnosingfixingcapability},
whereas gap-following changes it as teachers become easier or harder to imitate
\citep{sun2026d3mopd}. In either case, the same mixture controls what the
student learns and how much evidence estimates each task gradient.

A mixture is one knob doing two jobs. Sampling a task more often changes both
the objective the student minimizes and how many samples estimate that task's
gradient, so any mixture that adapts to training progress changes what is
learned and how precisely at the same time. The precision side is not a detail:
the noise of a task's gradient differs across tasks, drifts during training,
and reaches the optimizer only after it preconditions each coordinate, so the
task whose noise dominates a step depends on which optimizer is used.

On-policy distillation weights every sampled token by the difference between
the student's and the teacher's log-probability. Averaged under the student,
this difference is the local reverse-KL gap; after centering, its variation
across tokens interacts with the score function to determine the gradient.
At one position, two tokens of similar student probability can carry weights
of opposite sign and very different magnitude. That spread creates sampling
noise and can differ across teachers because teachers differ in how far they
are from the student. The teacher already computes a distribution at every
position, yet the usual sampled-token estimator keeps only one entry. We test
whether larger measured gaps predict larger spread rather than assuming that
relationship as a theorem.

We keep the objective fixed, compute each task's gradient with an estimator
that integrates the top-$k$ token contribution analytically, and spend
micro-batches on the noise that remains, measured in the coordinates the
optimizer's preconditioner defines.

\begin{figure}[t]
    \centering
    \includegraphics[width=\linewidth]{figures/gpas_main_figure.pdf}
    \caption{GPAS fixes the objective weights $w$ and adapts only the count
    vector $m=(m_1,\ldots,m_M)$. Each micro-batch gradient comes from the
    top-$k$ estimator. The completed step supplies the preconditioned gradient
    $Dg$, where $D$ comes from the optimizer and is the identity for SGD;
    measured time enables the optional cost-aware mode. No additional model
    pass is required. Segment widths and noise profiles are schematic.}
    \label{fig:gpas_overview}
\end{figure}

\missing{First empirical law: quantify the cross-task range and training drift
of preconditioned noise, and report how often its ranking differs from the
unpreconditioned ranking.} \missing{Second empirical law: report how the token
component tracks teacher--student gap and teacher origin, and how much the
top-$k$ estimator narrows the remaining cross-task spread.} \missing{Third
empirical law: at matched steps, report the separate and joint gains from
integrating the retained-token contribution and allocating the residual noise,
including the loss-gap control.} \missing{Fourth empirical law: report the overhead ratio and the GPU
hours saved by the cost-aware mode, together with every task's downstream
change relative to Uniform.}

Our contributions are:
\begin{itemize}
    \item a fixed-objective MOPD estimator whose micro-batch counts change
    variance but not expectation, with a local variance metric defined by any
    optimizer's pre-step diagonal preconditioner;
    \item a diagnosis connecting teacher--student gap to token-level noise, a
    conditionally unbiased top-$k$ estimator that integrates its retained-set
    component, and a complementary rule that allocates counts using the noise
    that remains;
    \item a four-teacher study of Qwen3-1.7B organized around the resulting
    empirical laws, with matched-step and matched-time comparisons and
    per-task capability measurements.
\end{itemize}
